% Propagazione degli errori nelle quattro operazioni
% OrbBits LaTeX conventions: see docs/style/latex.md
% Several languages: `orbbits build` defines \orbLocale per localised output; add
% \providecommand{\orbLocale}{it} here and switch babel and text on it (docs/localization.md).
\documentclass[11pt,a4paper]{article}
\usepackage[T1]{fontenc}
\usepackage[utf8]{inputenc}
\usepackage[italian]{babel}
\usepackage{amsmath,amssymb}
\usepackage[margin=2.2cm]{geometry}
\usepackage{graphicx}
\usepackage{xcolor}
\usepackage{enumitem}
\usepackage{microtype}
\usepackage[hidelinks]{hyperref}
\usepackage{booktabs}

\definecolor{orbPrimary}{HTML}{1F5FBF}

\title{Propagazione degli errori nelle quattro operazioni}
\author{}
\date{}

\begin{document}
\maketitle

\section*{Il problema}

I manuali enunciano le formule di propagazione dell'errore quasi sempre senza
dimostrarle. In questa dispensa le ricaviamo per esteso, con il metodo degli
\emph{estremi dell'intervallo di misura}: nessun passaggio saltato, e un
commento nei punti concettualmente delicati.

Partiamo da due grandezze misurate direttamente, nella forma consueta
\begin{equation*}
a = \bar a \pm \Delta a, \qquad b = \bar b \pm \Delta b,
\end{equation*}
dove $\bar a$ e $\bar b$ sono i valori misurati (valori nominali) e $\Delta a$,
$\Delta b$ gli errori assoluti. Il significato è: il valore vero di $a$ sta
nell'intervallo $[\bar a - \Delta a,\ \bar a + \Delta a]$, e lo stesso per $b$.
Conviene dare un nome agli estremi di questi intervalli:
\begin{equation*}
a_{\max} = \bar a + \Delta a, \qquad a_{\min} = \bar a - \Delta a,
\end{equation*}
\begin{equation*}
b_{\max} = \bar b + \Delta b, \qquad b_{\min} = \bar b - \Delta b.
\end{equation*}

Da $a$ e $b$ ricaviamo una grandezza \emph{indiretta} $c$, con una delle quattro
operazioni: $c = a+b$, $c = a-b$, $c = ab$, $c = a/b$. L'obiettivo è determinare
l'intervallo $[c_{\min},\, c_{\max}]$ in cui può trovarsi $c$ e riscriverlo
nella forma
\begin{equation*}
c = \bar c \pm \Delta c,
\end{equation*}
cioè trovare il valore nominale $\bar c$ e l'incertezza $\Delta c$ in funzione
di quelli di $a$ e $b$.

\paragraph{Errori massimi, non statistici.} Qui $\Delta a$ e $\Delta b$ sono
\emph{limiti di incertezza} (errori massimi): siamo cioè sicuri che il valore
vero cada dentro l'intervallo. Non trattiamo errori casuali indipendenti; per
questo nel seguito non comparirà mai la somma in quadratura
$\sqrt{\Delta a^2 + \Delta b^2}$, che appartiene a quel quadro diverso.

\paragraph{Ipotesi per prodotto e quoziente.} Nelle sezioni sul prodotto e sul
quoziente assumiamo esplicitamente
\begin{equation*}
\bar a > 0, \qquad \bar b > 0, \qquad \Delta a < \bar a, \qquad \Delta b < \bar b,
\end{equation*}
così che $a$ e $b$ restano positivi in tutto l'intervallo di misura; in
particolare l'intervallo di $b$ non contiene lo zero, condizione indispensabile
per poter dividere. Queste ipotesi evitano problemi di monotonia: garantiscono
che, dentro gli intervalli, le funzioni che useremo crescano o decrescano sempre
nello stesso verso, e quindi che massimo e minimo si ottengano agli estremi.

\section*{La somma: $c = a + b$}

Il valore nominale della somma è
\begin{equation*}
\bar c = \bar a + \bar b.
\end{equation*}
La somma cresce sia al crescere di $a$ sia al crescere di $b$: il massimo di $c$
si ottiene prendendo entrambe le grandezze al massimo, il minimo prendendole
entrambe al minimo. Svolgiamo i conti per il massimo:
\begin{align*}
c_{\max} &= a_{\max} + b_{\max} && \text{entrambe al massimo}\\
&= (\bar a + \Delta a) + (\bar b + \Delta b) && \text{per definizione di estremi}\\
&= (\bar a + \bar b) + (\Delta a + \Delta b) && \text{raggruppando}\\
&= \bar c + (\Delta a + \Delta b) && \text{per definizione di } \bar c
\end{align*}
e per il minimo, in modo del tutto analogo:
\begin{align*}
c_{\min} &= a_{\min} + b_{\min} && \text{entrambe al minimo}\\
&= (\bar a - \Delta a) + (\bar b - \Delta b)\\
&= (\bar a + \bar b) - (\Delta a + \Delta b)\\
&= \bar c - (\Delta a + \Delta b).
\end{align*}
L'intervallo di $c$ è quindi esattamente simmetrico rispetto a $\bar c$, e
l'incertezza della somma è
\begin{equation*}
\boxed{\,\Delta c = \Delta a + \Delta b\,}
\tag*{esatta}
\end{equation*}
sicché la misura indiretta si scrive
\begin{equation*}
\boxed{\,c = (\bar a + \bar b) \pm (\Delta a + \Delta b)\,}
\tag*{esatta}
\end{equation*}

\paragraph{Osservazione.} Nel modello adottato (incertezze come limiti) questo
risultato non è un'approssimazione: è la conseguenza esatta delle definizioni.
Nella somma, gli errori assoluti si sommano.

\section*{La differenza: $c = a - b$}

Ora il valore nominale è
\begin{equation*}
\bar c = \bar a - \bar b,
\end{equation*}
e il passaggio da capire bene riguarda \emph{quali} estremi accoppiare.

\paragraph{Massimo e minimo si accoppiano in modo incrociato.} Al crescere di
$a$ la differenza $a-b$ cresce; al crescere di $b$ essa \emph{diminuisce}.
Quindi $c$ è massima quando $a$ è al massimo e $b$ è al minimo, ed è minima
quando $a$ è al minimo e $b$ è al massimo. In parole: un errore che fa sembrare
$b$ più grande di quanto sia rende $a-b$ più piccolo; ecco perché gli estremi
vanno accoppiati in modo incrociato.

Svolgiamo il massimo:
\begin{align*}
c_{\max} &= a_{\max} - b_{\min} && \text{estremi incrociati}\\
&= (\bar a + \Delta a) - (\bar b - \Delta b) && \text{per definizione di estremi}\\
&= \bar a + \Delta a - \bar b + \Delta b && \text{attenzione ai segni}\\
&= (\bar a - \bar b) + (\Delta a + \Delta b) && \text{raggruppando}\\
&= \bar c + (\Delta a + \Delta b)
\end{align*}
Soffermiamoci sul passo segnalato: il segno meno davanti alla parentesi cambia
il segno a \emph{entrambi} i termini, $-(\bar b - \Delta b) = -\bar b +
\Delta b$. È proprio qui che il meno della differenza, agendo su $-\Delta b$,
produce un più: l'incertezza di $b$ entra con il segno positivo.

Per il minimo:
\begin{align*}
c_{\min} &= a_{\min} - b_{\max} && \text{estremi incrociati}\\
&= (\bar a - \Delta a) - (\bar b + \Delta b)\\
&= \bar a - \Delta a - \bar b - \Delta b && \text{stessa attenzione ai segni}\\
&= (\bar a - \bar b) - (\Delta a + \Delta b)\\
&= \bar c - (\Delta a + \Delta b).
\end{align*}
Anche qui l'intervallo è esattamente simmetrico rispetto a $\bar c$, quindi
\begin{equation*}
\boxed{\,\Delta c = \Delta a + \Delta b\,}
\tag*{esatta}
\end{equation*}
e
\begin{equation*}
\boxed{\,c = (\bar a - \bar b) \pm (\Delta a + \Delta b)\,}
\tag*{esatta}
\end{equation*}

\paragraph{Perché gli errori si sommano anche sottraendo.} Il segno meno della
formula riguarda i valori nominali, non le incertezze: sottrarre due misure non
\emph{compensa} gli errori, li accumula, perché il caso peggiore è quello in cui
i due errori spingono in versi opposti. Vale la pena notare una conseguenza:
sottraendo grandezze quasi uguali ($\bar a \approx \bar b$) il valore nominale
$\bar c$ diventa piccolo mentre $\Delta c$ resta la somma delle incertezze, e
l'errore \emph{relativo} del risultato può diventare molto grande anche partendo
da misure ottime.

\section*{Il prodotto: $c = ab$}

Vale l'ipotesi fatta all'inizio: $\bar a > 0$, $\bar b > 0$, $\Delta a < \bar a$,
$\Delta b < \bar b$. I due fattori restano quindi positivi in tutto l'intervallo
e il prodotto cresce al crescere di ciascuno di essi: massimo e minimo si
ottengono accoppiando gli estremi nello stesso verso. Il valore nominale è
\begin{equation*}
\bar c = \bar a\,\bar b.
\end{equation*}

\paragraph{Massimo e minimo esatti.} Per il massimo:
\begin{align*}
c_{\max} &= a_{\max}\, b_{\max}
= (\bar a + \Delta a)(\bar b + \Delta b)\\
&= \bar a\,\bar b + \bar a\,\Delta b + \bar b\,\Delta a + \Delta a\,\Delta b
&& \text{sviluppo completo}
\end{align*}
e quindi, essendo $\bar c = \bar a\,\bar b$,
\begin{equation*}
c_{\max} - \bar c
= \bar a\,\Delta b + \bar b\,\Delta a + \Delta a\,\Delta b .
\end{equation*}
Per il minimo:
\begin{align*}
c_{\min} &= a_{\min}\, b_{\min}
= (\bar a - \Delta a)(\bar b - \Delta b)\\
&= \bar a\,\bar b - \bar a\,\Delta b - \bar b\,\Delta a + \Delta a\,\Delta b
&& \text{nota: } (-\Delta a)(-\Delta b) = +\Delta a\,\Delta b
\end{align*}
da cui
\begin{equation*}
\bar c - c_{\min}
= \bar a\,\Delta b + \bar b\,\Delta a - \Delta a\,\Delta b .
\end{equation*}

\paragraph{L'intervallo esatto non è simmetrico.} Lo scarto superiore
$c_{\max} - \bar c$ e quello inferiore $\bar c - c_{\min}$ differiscono per il
termine $\Delta a\,\Delta b$, che compare con $+$ nel primo e con $-$ nel
secondo: rispetto a $\bar c$ l'intervallo esatto è leggermente sbilanciato verso
l'alto. Il colpevole è proprio il termine $\Delta a\,\Delta b$, che nelle somme
e nelle differenze non era mai comparso.

\paragraph{Errori relativi.} Conviene introdurre gli errori relativi
\begin{equation*}
\varepsilon_a = \frac{\Delta a}{\bar a},
\qquad
\varepsilon_b = \frac{\Delta b}{\bar b}
\end{equation*}
e dividere gli scarti per $\bar c = \bar a\,\bar b$. Per l'estremo superiore:
\begin{align*}
\frac{c_{\max} - \bar c}{\bar c}
&= \frac{\bar a\,\Delta b + \bar b\,\Delta a + \Delta a\,\Delta b}{\bar a\,\bar b}\\
&= \frac{\Delta b}{\bar b} + \frac{\Delta a}{\bar a}
+ \frac{\Delta a}{\bar a}\cdot\frac{\Delta b}{\bar b}
&& \text{termine a termine}\\
&= \varepsilon_a + \varepsilon_b + \varepsilon_a\,\varepsilon_b ,
\end{align*}
e, con lo stesso conto per il minimo,
\begin{equation*}
\frac{\bar c - c_{\min}}{\bar c}
= \varepsilon_a + \varepsilon_b - \varepsilon_a\,\varepsilon_b .
\end{equation*}

\paragraph{Il termine $\varepsilon_a\,\varepsilon_b$ è di secondo ordine.} Se
gli errori relativi sono piccoli --- in un laboratorio scolastico tipicamente
dell'ordine dell'1\% o meno --- il loro prodotto è cento volte più piccolo di
ciascuno di essi: ad esempio $\varepsilon_a = \varepsilon_b = 1\%$ dà
$\varepsilon_a\,\varepsilon_b = 0{,}01\%$. Un termine che è il prodotto di due
piccole quantità si dice \emph{di secondo ordine} nelle incertezze relative.
Trascurandolo, i due scarti coincidono e l'intervallo diventa, con ottima
approssimazione, simmetrico: $c = \bar c \pm \Delta c$.

\paragraph{Come leggere le formule che seguono.} Da qui in poi etichettiamo i
risultati in tre modi: \emph{esatta} (conseguenza esatta delle definizioni),
\emph{al primo ordine} (ottenuta trascurando i termini di secondo ordine nelle
incertezze relative) e \emph{formula d'uso} (la forma sotto cui la si incontra
nei problemi scolastici, ed è la stessa approssimazione al primo ordine
riscritta con gli errori assoluti).

Trascurando $\varepsilon_a\,\varepsilon_b$ otteniamo la formula usuale per
l'errore relativo del prodotto:
\begin{equation*}
\boxed{\,\frac{\Delta c}{\bar c} \simeq \frac{\Delta a}{\bar a} + \frac{\Delta b}{\bar b}\,}
\qquad\text{ovvero}\qquad
\boxed{\,\varepsilon_c \simeq \varepsilon_a + \varepsilon_b\,}
\tag*{al primo ordine}
\end{equation*}

\paragraph{La forma con gli errori assoluti.} Per l'errore assoluto basta
moltiplicare per $\bar c = \bar a\,\bar b$ e semplificare termine a termine:
\begin{align*}
\Delta c
&\simeq \bar a\,\bar b \left( \frac{\Delta a}{\bar a} + \frac{\Delta b}{\bar b} \right)\\
&= \bar a\,\bar b \cdot \frac{\Delta a}{\bar a}
+ \bar a\,\bar b \cdot \frac{\Delta b}{\bar b}
&& \text{termine a termine}\\
&= \bar b\,\Delta a + \bar a\,\Delta b
\end{align*}
cioè
\begin{equation*}
\boxed{\,\Delta c \simeq \bar b\,\Delta a + \bar a\,\Delta b\,}
\tag*{formula d'uso}
\end{equation*}

\section*{Il quoziente: $c = a/b$}

Anche qui valgono le ipotesi iniziali; in particolare $b_{\min} = \bar b -
\Delta b > 0$, quindi il denominatore non cambia mai segno e possiamo dividere.
Il valore nominale è
\begin{equation*}
\bar c = \frac{\bar a}{\bar b}.
\end{equation*}

\paragraph{Quali estremi producono massimo e minimo.} Il rapporto cresce al
crescere del numeratore e --- restando il denominatore positivo --- al
\emph{decrescere} del denominatore: aumentando il numeratore il rapporto sale,
aumentando il denominatore scende. Quindi
\begin{equation*}
c_{\max} = \frac{a_{\max}}{b_{\min}} = \frac{\bar a + \Delta a}{\bar b - \Delta b},
\qquad
c_{\min} = \frac{a_{\min}}{b_{\max}} = \frac{\bar a - \Delta a}{\bar b + \Delta b}:
\end{equation*}
di nuovo estremi accoppiati in modo incrociato.

Riscriviamo $c_{\max}$ in termini degli errori relativi, raccogliendo $\bar a$ e
$\bar b$:
\begin{equation*}
c_{\max}
= \frac{\bar a\,(1 + \varepsilon_a)}{\bar b\,(1 - \varepsilon_b)}
= \bar c\,\frac{1 + \varepsilon_a}{1 - \varepsilon_b}.
\end{equation*}
Calcoliamo \emph{esattamente} lo scarto relativo superiore:
\begin{align*}
\frac{c_{\max} - \bar c}{\bar c}
&= \frac{1 + \varepsilon_a}{1 - \varepsilon_b} - 1\\
&= \frac{1 + \varepsilon_a - (1 - \varepsilon_b)}{1 - \varepsilon_b}
&& \text{al denominatore comune}\\
&= \frac{\varepsilon_a + \varepsilon_b}{1 - \varepsilon_b},
\end{align*}
\begin{equation*}
\boxed{\,\frac{c_{\max} - \bar c}{\bar c} = \frac{\varepsilon_a + \varepsilon_b}{1 - \varepsilon_b}\,}
\tag*{esatta}
\end{equation*}
Analogamente, partendo da
\begin{equation*}
c_{\min} = \bar c\,\frac{1 - \varepsilon_a}{1 + \varepsilon_b},
\end{equation*}
lo scarto relativo inferiore è
\begin{align*}
\frac{\bar c - c_{\min}}{\bar c}
&= 1 - \frac{1 - \varepsilon_a}{1 + \varepsilon_b}
= \frac{(1 + \varepsilon_b) - (1 - \varepsilon_a)}{1 + \varepsilon_b}
= \frac{\varepsilon_a + \varepsilon_b}{1 + \varepsilon_b},
\end{align*}
\begin{equation*}
\boxed{\,\frac{\bar c - c_{\min}}{\bar c} = \frac{\varepsilon_a + \varepsilon_b}{1 + \varepsilon_b}\,}
\tag*{esatta}
\end{equation*}

\paragraph{Anche qui, asimmetria.} I due scarti non sono esattamente uguali:
il primo ha denominatore $1 - \varepsilon_b < 1$, il secondo $1 + \varepsilon_b
> 1$, quindi lo scarto superiore è un poco più grande di quello inferiore. È lo
stesso fenomeno del termine $\Delta a\,\Delta b$ nel prodotto, vestito in modo
diverso.

\paragraph{L'approssimazione.} Per errori relativi piccoli, $\varepsilon_b \ll
1$, il denominatore $1 - \varepsilon_b$ è molto vicino a $1$; più precisamente
si può scrivere
\begin{equation*}
\frac{1}{1 - \varepsilon_b} \simeq 1 + \varepsilon_b,
\end{equation*}
giacché $(1 - \varepsilon_b)(1 + \varepsilon_b) = 1 - \varepsilon_b^2 \simeq 1$.
Sostituendo, comparirebbe il termine $(\varepsilon_a + \varepsilon_b)\,
\varepsilon_b$, che è di secondo ordine e si trascura come nel prodotto.
Restano così le formule note:
\begin{equation*}
\boxed{\,\frac{\Delta c}{\bar c} \simeq \frac{\Delta a}{\bar a} + \frac{\Delta b}{\bar b}\,}
\qquad\text{ovvero}\qquad
\boxed{\,\varepsilon_c \simeq \varepsilon_a + \varepsilon_b\,}
\tag*{al primo ordine}
\end{equation*}

\paragraph{La forma con gli errori assoluti.} Moltiplicando per $\bar c =
\bar a / \bar b$ e sviluppando i due termini:
\begin{align*}
\Delta c
&\simeq \frac{\bar a}{\bar b} \left( \frac{\Delta a}{\bar a} + \frac{\Delta b}{\bar b} \right)\\
&= \frac{\bar a}{\bar b} \cdot \frac{\Delta a}{\bar a}
+ \frac{\bar a}{\bar b} \cdot \frac{\Delta b}{\bar b}
&& \text{termine a termine}\\
&= \frac{\Delta a}{\bar b} + \frac{\bar a}{\bar b^{2}}\,\Delta b,
\end{align*}
\begin{equation*}
\boxed{\,\Delta c \simeq \frac{\Delta a}{\bar b} + \frac{\bar a}{\bar b^{2}}\,\Delta b\,}
\tag*{formula d'uso}
\end{equation*}

\section*{Le quattro formule a confronto}

{\small
\begin{center}
\begin{tabular}{@{}c l l@{}}
\toprule
Grandezza & Incertezze di $c$ & Natura del risultato \\
\midrule
$a+b$ & $\Delta c = \Delta a + \Delta b$ & esatta \\[3pt]
$a-b$ & $\Delta c = \Delta a + \Delta b$ & esatta \\[3pt]
$ab$ &
\begin{tabular}[t]{@{}l@{}}
$\dfrac{\Delta c}{\bar c} \simeq \dfrac{\Delta a}{\bar a} + \dfrac{\Delta b}{\bar b}$\\[5pt]
$\Delta c \simeq \bar b\,\Delta a + \bar a\,\Delta b$
\end{tabular}
& approssimata al primo ordine \\[3pt]
$\dfrac{a}{b}$ &
\begin{tabular}[t]{@{}l@{}}
$\dfrac{\Delta c}{\bar c} \simeq \dfrac{\Delta a}{\bar a} + \dfrac{\Delta b}{\bar b}$\\[5pt]
$\Delta c \simeq \dfrac{\Delta a}{\bar b} + \dfrac{\bar a}{\bar b^{2}}\,\Delta b$
\end{tabular}
& approssimata al primo ordine \\
\bottomrule
\end{tabular}
\end{center}
}

\paragraph{Nota finale.} Le formule per somma e differenza sono conseguenze
esatte delle definizioni. Le formule per prodotto e quoziente sono invece
approssimate \emph{al primo ordine}: valgono quando le incertezze relative sono
piccole, e il passaggio che le produce --- trascurare $\varepsilon_a
\varepsilon_b$ e sostituire $1/(1 - \varepsilon_b)$ con $1 + \varepsilon_b$ ---
è un'approssimazione reale, non una semplificazione cosmetica. Se servisse un
valore fedele o un limite garantito, le formule esatte sono quelle ricavate
negli scarti di prodotto e quoziente. Con incertezze dell'ordine di pochi
punti percentuale --- il caso normale nei problemi scolastici --- l'errore
commesso usando le formule al primo ordine è molto più piccolo delle
incertezze stesse, ed è per questo che esse sono le formule standard.

\end{document}
